%%%%%%%%%%%%%%%%%%%%%%%%%%%%%%%%%%%%%%%%%%%%%%%%%%%%%%%%%%%%%%%%%%%%%%%%%%%%%%%
%
%  ALGORITHMS & EQUATIONS — INSERT GUIDE
%  Paper: "Beyond the Baseline: Improving LLM-driven Agents for ACD"
%
%  This file contains all formal definitions, equations, and algorithms
%  extracted from the four related papers, adapted for YOUR paper.
%  Each block is labeled with the EXACT section where it should be inserted.
%
%  REQUIRED: Add these packages to your preamble (if not already present):
%    \usepackage{amsmath}
%    \usepackage{amssymb}
%    \usepackage{algorithm}
%    \usepackage{algorithmic}
%
%%%%%%%%%%%%%%%%%%%%%%%%%%%%%%%%%%%%%%%%%%%%%%%%%%%%%%%%%%%%%%%%%%%%%%%%%%%%%%%


%=============================================================================
%  SECTION 2: BACKGROUND AND TECHNICAL FOUNDATIONS
%=============================================================================
%
%  INSERT INTO: Section 2.1 — "The CybORG and CAGE 4 Environment"
%  PURPOSE:     Formalize the Dec-POMDP that underpins the ENTIRE paper.
%               Without this, the paper lacks a mathematical problem statement.
%  SOURCE:      Singh et al. (2025) Section 3 + King et al. (2025) Section 3
%-----------------------------------------------------------------------------

% --- Equation: Dec-POMDP Definition ---
% INSERT after the sentence: "...meaning each defender only has a limited 
% view of the overall network state."

The CAGE~4 environment is formally modeled as a
\textbf{Decentralized Partially Observable Markov Decision Process}
(Dec-POMDP)~\cite{arXiv:2410.17351}, defined by the tuple:
\begin{equation}
\label{eq:dec_pomdp}
  \mathcal{M} = \bigl(\mathcal{I},\; \mathcal{S},\; \mathcal{A},\;
  \mathcal{T},\; \Omega,\; O,\; R,\; b_0 \bigr),
\end{equation}
where $\mathcal{I}=\{1,\ldots,n\}$ is the set of $n$ blue agents,
$\mathcal{S}$ is a finite state space capturing the full network state,
$\mathcal{A}=\times_{i\in\mathcal{I}} \mathcal{A}_i$ is the
joint action space composed of each agent's individual actions
(Monitor, Analyse, Remove, Restore, DeployDecoy,
AllowTrafficZone, BlockTrafficZone),
$\mathcal{T}: \mathcal{S}\times\mathcal{A}\to\Delta(\mathcal{S})$
is the transition function,
$\Omega=\times_{i\in\mathcal{I}}\Omega_i$ is the set of
joint observations,
$O:\mathcal{S}\times\mathcal{A}\to\Omega$ is the observation function,
$R:\mathcal{S}\times\mathcal{A}\to\mathbb{R}$ is the shared
reward function (penalties for service unavailability),
and $b_0$ is the initial state distribution.
Each blue agent $i$ receives only its local observation
$o_i\in\Omega_i$ and must select actions based on its observation
history $h_t^i = (o_0^i, a_0^i, \ldots, o_{t-1}^i, a_{t-1}^i)$.


%=============================================================================
%  SECTION 2: BACKGROUND AND TECHNICAL FOUNDATIONS
%=============================================================================
%
%  INSERT INTO: Section 2.2 — "The KEEP Agent Baseline"
%  PURPOSE:     Show the GCN equation that gives KEEP its power.
%               This explains WHY KEEP outperforms flat-vector LLM prompts.
%  SOURCE:      King et al. (2025) Section 4.2, Equation 2
%-----------------------------------------------------------------------------

% --- Equation: GCN Node Embedding ---
% INSERT after: "Developed using the Proximal Policy Optimization (PPO) 
% algorithm..."

The KEEP architecture represents the network as an attributed graph
$\mathcal{O}=\langle V, E, X \rangle$, where $V$ denotes hosts,
$E$ denotes inter-host connections, and $X$ denotes per-node feature
vectors. A two-layer \textbf{Graph Convolutional Network}
(GCN)~\cite{arXiv:2509.16151} produces node embeddings via:
\begin{equation}
\label{eq:gcn}
  H^{(\ell+1)} = \sigma\!\left(
    \tilde{D}^{-\frac{1}{2}}\,\tilde{A}\,\tilde{D}^{-\frac{1}{2}}\,
    H^{(\ell)}\,W^{(\ell)} + b^{(\ell)}
  \right),
\end{equation}
where $\tilde{A}=A+I$ is the adjacency matrix with self-loops,
$\tilde{D}$ is its degree matrix,
$W^{(\ell)}$ and $b^{(\ell)}$ are trainable parameters,
and $\sigma$ is a non-linear activation function.
The first GCN layer uses 256-dimensional embeddings and the second
uses 64-dimensional embeddings.
Actions are represented as functions upon nodes $a_k(v_i)$,
allowing the action space to scale automatically with the network
topology. This relational inductive bias enables KEEP to
zero-shot generalize across unseen topologies --- an advantage
that explains its strong performance ($\mu=-493$) relative to
agents operating on flat observation vectors.


%=============================================================================
%  SECTION 2: BACKGROUND AND TECHNICAL FOUNDATIONS
%=============================================================================
%
%  INSERT INTO: Section 2.2 — "The KEEP Agent Baseline"
%  PURPOSE:     Show the PPO clipped objective that trains KEEP.
%               Needed because PPO is referenced throughout the paper
%               but never formally defined.
%  SOURCE:      Schulman et al. (2017) via Singh et al. Section 4
%-----------------------------------------------------------------------------

% --- Equation: PPO Clipped Objective ---

Both KEEP and the H-MARL sub-policies are trained using
\textbf{Proximal Policy Optimization} (PPO)~\cite{arXiv:1707.06347},
which maximizes the clipped surrogate objective:
\begin{equation}
\label{eq:ppo}
  L^{\text{CLIP}}(\theta) = \hat{\mathbb{E}}_t\!\left[
    \min\!\bigl(
      r_t(\theta)\,\hat{A}_t,\;
      \text{clip}\bigl(r_t(\theta),\, 1{-}\epsilon,\, 1{+}\epsilon\bigr)\,
      \hat{A}_t
    \bigr)
  \right],
\end{equation}
where $r_t(\theta) = \frac{\pi_\theta(a_t|s_t)}{\pi_{\theta_{\text{old}}}(a_t|s_t)}$
is the probability ratio between the new and old policies,
$\hat{A}_t$ is the generalized advantage estimate (GAE),
and $\epsilon$ is a clipping hyperparameter (typically $0.2$).
This objective prevents destructive large policy updates,
enabling stable multi-agent training.


%=============================================================================
%  SECTION 3: RELATED WORK — Hammar et al. subsection
%=============================================================================
%
%  INSERT INTO: Section 3.3 — "Hallucination-Resistant Security Planning"
%  PURPOSE:     Define eq:consistency (ALREADY REFERENCED in your Table I
%               at line 346 but NEVER DEFINED — this is a broken reference).
%  SOURCE:      Hammar et al. (2026) Equation 2
%-----------------------------------------------------------------------------

% --- Equation: Lookahead Consistency Function ---
% INSERT after: "...checking how consistent they are with each other."

Formally, given $N$ candidate actions
$\mathcal{A}_t = \{a^1_t, \ldots, a^N_t\}$,
each is paired with a lookahead prediction $T^i_{t+1}$ of the
expected time remaining to complete the task.
The \textbf{consistency} of the candidate set is:
\begin{equation}
\label{eq:consistency}
  \lambda(\mathcal{A}_t) = \exp\!\left(
    -\frac{\beta}{N}\sum_{i=1}^{N}
    \bigl(T^i_{t+1} - \overline{T}_{t+1}\bigr)^2
  \right),
\end{equation}
where $\beta > 0$ controls sensitivity to disagreement and
$\overline{T}_{t+1} = \frac{1}{N}\sum_{i=1}^{N} T^i_{t+1}$
is the mean prediction.
The output is bounded in $[0,1]$: a value of~$1$ indicates
full consistency (all candidates agree) and~$0$ indicates
high disagreement, signaling likely hallucination.


%=============================================================================
%  SECTION 3: RELATED WORK — Hammar et al. subsection
%=============================================================================
%
%  INSERT INTO: Section 3.3 — immediately after eq:consistency
%  PURPOSE:     The abstention policy that uses the consistency score.
%               This is the mechanism that inspired your CoSC fallback.
%  SOURCE:      Hammar et al. (2026) Equation 3
%-----------------------------------------------------------------------------

% --- Equation: Conformal Abstention Policy ---

The framework selects actions via the \textbf{conformal abstention
policy}:
\begin{equation}
\label{eq:abstention}
  \pi_\gamma(\mathcal{A}_t) =
  \begin{cases}
    \emptyset \;\text{(abstain)},
      & \text{if } \lambda(\mathcal{A}_t) \leq \gamma, \\[4pt]
    \displaystyle\arg\min_{a^i_t \in \mathcal{A}_t}\{T^i_{t+1}\},
      & \text{if } \lambda(\mathcal{A}_t) > \gamma,
  \end{cases}
\end{equation}
where $\gamma \in [0,1]$ is a consistency threshold.
When consistency falls below~$\gamma$, the system abstains and
collects external feedback (e.g., from a digital twin) before
retrying. Hammar et al.\ prove that tuning~$\gamma$ on a
calibration dataset of $n$ known-hallucinated examples provides a
formal upper bound on the hallucination probability: given a
desired bound $\kappa\in(0,1]$, the threshold is set to
\begin{equation}
\label{eq:gamma_calibration}
  \tilde{\gamma} = \inf\!\left\{
    \gamma \;\middle|\;
    \frac{|\{i \mid \lambda(\mathcal{A}_i)\leq\gamma\}|}{n}
    \geq
    \frac{\lceil(n+1)(1-\kappa)\rceil}{n}
  \right\},
\end{equation}
guaranteeing $P\!\bigl(\pi_{\tilde{\gamma}}(\tilde{\mathcal{A}})
\neq\emptyset\bigr)\leq\kappa$\;\cite{arXiv:2602.05279}.


%=============================================================================
%  SECTION 3: RELATED WORK — Hammar et al. subsection
%=============================================================================
%
%  INSERT INTO: Section 3.3 — near "hallucination rate of 0.02"
%  PURPOSE:     Formal definition of a hallucinated action.
%               Your paper measures hallucination rate but never
%               formally defines what counts as a hallucination.
%  SOURCE:      Hammar et al. (2026) Definition 1
%-----------------------------------------------------------------------------

% --- Definition: Hallucinated Action ---

\noindent\textbf{Definition 1 (Hallucinated Action~\cite{arXiv:2602.05279}).}
Given a sequence of actions $a_0, \ldots, a_{t-1}, a_t$ for a
security management task, the action~$a_t$ is \textit{hallucinated}
if it does not reduce the expected time remaining to complete the task:
\begin{equation}
\label{eq:hallucination_def}
  \mathbb{E}\{T_{t+1} \mid a_0, \ldots, a_t\}
  \;\geq\;
  \mathbb{E}\{T_t \mid a_0, \ldots, a_{t-1}\},
\end{equation}
where $T_t$ is the time remaining at stage~$t$, assuming future
actions are selected optimally.


%=============================================================================
%  SECTION 3: RELATED WORK — Gao et al. subsection
%=============================================================================
%
%  INSERT INTO: Section 3.4 — "In-Context Autonomous Network Incident
%               Response"
%  PURPOSE:     The LoRA fine-tuning loss and Monte-Carlo Q-function.
%               These are the two key equations from Gao et al. that
%               directly inform your CoT template and future work on
%               behavioral distillation.
%  SOURCE:      Gao et al. (2026) Equations 2 and 3
%-----------------------------------------------------------------------------

% --- Equation: LoRA Fine-Tuning Cross-Entropy Loss ---
% INSERT after: "The model is fine-tuned with LoRA on 50,000 labeled 
% examples"

The model is fine-tuned offline using \textbf{Low-Rank Adaptation}
(LoRA), minimizing the cross-entropy loss over a dataset
$\mathcal{D}=\{(x^i, y^i)\}_{i=1}^K$ of instruction--answer pairs:
\begin{equation}
\label{eq:lora_loss}
  L(w) = -\frac{1}{B}\sum_{i=1}^{B}\sum_{k=1}^{\ell_i}
  \log\,\Phi_w\!\bigl(y^i_k \mid x^i, y^i_{1:k-1}\bigr),
\end{equation}
where $w$ denotes the LoRA trainable parameters (rank~64),
$B$ is the batch size, and $\ell_i$ is the token length of
the $i$-th answer.


% --- Equation: Monte-Carlo Lookahead Q-Function ---
% INSERT after eq:lora_loss, in the planning description

Online, the agent evaluates $N$ candidate actions by simulating $M$
recovery trajectories for each, then selects the cost minimizer
via a Monte-Carlo estimate of the Q-function:
\begin{equation}
\label{eq:mc_q}
  Q(\hat{s}_{t+1}, \hat{a}^k_{t+1})
  = \frac{1}{M}\sum_{i\in[M]}\sum_{(\hat{s},\hat{a})\in q^i}
    c(\hat{s}, \hat{a}),
\end{equation}
\begin{equation}
\label{eq:mc_argmin}
  a_{t+1} \in \arg\min_{a \in \mathcal{A}_t}\;
  Q(\hat{s}_{t+1},\, a),
\end{equation}
where $q^i$ is the $i$-th simulated trajectory and $c(\hat{s},\hat{a})$
is the time cost of executing action~$\hat{a}$ in state~$\hat{s}$.
This planning procedure has complexity $O(MN)$ per step.


%=============================================================================
%  SECTION 4: METHODOLOGY
%=============================================================================
%
%  INSERT INTO: Section 4.1 — "Agent Architecture and CoSC Extension"
%  PURPOSE:     The IOC-Priority Expert Rule that forms the semantic
%               verification logic INSIDE your CoSC.  This is the
%               "reference for what correct reasoning looks like" that
%               you mention in the Related Work but never formalize.
%  SOURCE:      Singh et al. (2025) Section 4.1, H-MARL Expert Policy
%-----------------------------------------------------------------------------

% --- Equation: IOC-Based Expert Decision Rule ---
% INSERT after the CoSC description bullet point

The semantic verification layer within the CoSC gate is grounded
in the \textbf{IOC-based expert policy}
$\pi_E$~\cite{arXiv:2410.17351}, which defines the
domain-correct reasoning reference:
\begin{equation}
\label{eq:ioc_expert}
  \pi_E(h_t) =
  \begin{cases}
    \texttt{Recover},
      & \text{if malicious-file IOCs detected on host}, \\
    \texttt{ControlTraffic},
      & \text{if network IOCs detected}, \\
    \texttt{Investigate},
      & \text{otherwise}.
  \end{cases}
\end{equation}
In our implementation, this rule is encoded as a
\textit{post-generation validation check}: if the LLM proposes
an action that contradicts $\pi_E$ given the current IOC state
(e.g., issuing \texttt{Analyse} when host-level IOCs indicate
active compromise), the CoSC gate flags a semantic
hallucination and triggers re-generation with corrective
feedback.


%=============================================================================
%  SECTION 4: METHODOLOGY
%=============================================================================
%
%  INSERT INTO: Section 4.1 — REPLACE the prose-only CoSC description
%               with this formal algorithm.
%  PURPOSE:     This is your CORE CONTRIBUTION formalized as pseudocode.
%               An academic paper must present its novel mechanism as
%               a reproducible algorithm.
%  SOURCE:      YOUR novel synthesis of:
%               - Hammar et al.'s consistency/abstention framework
%               - Singh et al.'s IOC expert rule
%               - Gao et al.'s perceive-reason-plan-act pipeline
%               - Castro et al.'s adapter interface
%-----------------------------------------------------------------------------

% --- Algorithm 1: Chain-of-Self-Correction (CoSC) ---

\begin{algorithm}[t]
\caption{Chain-of-Self-Correction (CoSC) Agent Policy}
\label{alg:cosc}
\begin{algorithmic}[1]
\REQUIRE LLM model $\Phi$, CybORG adapter $\mathcal{F}$,
         IOC expert rule $\pi_E$, max retries $K=3$
\ENSURE  Action $a_t$ for current step

\STATE $o_t \gets$ CybORG observation at step $t$
\STATE $\textit{prompt} \gets \mathcal{F}(o_t)$
    \COMMENT{Adapter: numerical obs $\to$ structured NL}
\STATE $\textit{ioc\_state} \gets$ extract IOC priorities from $o_t$
\STATE $k \gets 0$

\WHILE{$k < K$}
    \STATE $(\hat{a}_t,\, \textit{reason}) \gets \Phi(\textit{prompt})$
        \COMMENT{LLM generates candidate action + reasoning}

    \STATE \textbf{// Stage 1: Schema Validation (Syntactic Check)}
    \IF{$\hat{a}_t \notin \mathcal{A}_{\text{valid}}$ \textbf{or}
        JSON parse fails}
        \STATE $\textit{prompt} \gets \textit{prompt} \,\cup\,
               \text{``Invalid format. Valid actions: } \mathcal{A}_{\text{valid}}\text{''}$
        \STATE $k \gets k + 1$;\; \textbf{continue}
    \ENDIF

    \STATE \textbf{// Stage 2: IOC-Rule Verification (Semantic Check)}
    \STATE $a_E \gets \pi_E(\textit{ioc\_state})$
        \COMMENT{Expert reference action (Eq.~\ref{eq:ioc_expert})}
    \IF{$\textit{ioc\_state}$ indicates active compromise
        \textbf{and} $\hat{a}_t \notin \{\texttt{Remove}, \texttt{Restore}\}$}
        \STATE $\textit{prompt} \gets \textit{prompt} \,\cup\,
               \text{``IOCs detected on host. Re-assess: recovery action required.''}$
        \STATE $k \gets k + 1$;\; \textbf{continue}
    \ENDIF

    \STATE \textbf{// Stage 3: Logical Consistency Check}
    \IF{$\hat{a}_t = \texttt{Restore}$ \textbf{and}
        target host has $\textit{ioc\_priority}=0$}
        \STATE $\textit{prompt} \gets \textit{prompt} \,\cup\,
               \text{``Host is clean. Restore would cause unnecessary downtime.''}$
        \STATE $k \gets k + 1$;\; \textbf{continue}
    \ENDIF

    \RETURN $\hat{a}_t$
        \COMMENT{Action passes all gates}
\ENDWHILE

\RETURN $\texttt{Sleep}$
    \COMMENT{Exhausted retries $\to$ safe fallback (abstention)}
\end{algorithmic}
\end{algorithm}


%=============================================================================
%  SECTION 4: METHODOLOGY — Evaluation Metrics
%=============================================================================
%
%  INSERT INTO: Section 4.4 — "Evaluation Metrics"
%  PURPOSE:     This is the explicit TODO on line 424 of your paper:
%               "demonstrate how the calculate was done - show the equations"
%  SOURCE:      Castro et al. (reward), Singh et al. (MTTR, precision),
%               Hammar et al. (hallucination), your own Cv definition
%-----------------------------------------------------------------------------

% --- REPLACE the current Section 4.4 content with the following ---

\subsection{Evaluation Metrics}
We quantify agent effectiveness using metrics established across
the reviewed literature. All metrics are computed over $E$ episodes
of $T$ steps each, with $n$ blue agents per episode.

\textbf{Joint Reward} ($\mu$).
The primary measure of team efficacy, computed as the mean
cumulative penalty across episodes. CybORG assigns penalties
when green agents cannot access services
(due to defensive actions or red-agent impact):
\begin{equation}
\label{eq:joint_reward}
  \mu = \frac{1}{E}\sum_{e=1}^{E}\sum_{t=1}^{T} R_t^{(e)},
\end{equation}
where $R_t^{(e)}$ is the shared team reward at step~$t$ of
episode~$e$~\cite{arXiv:2505.04843v2}.

\textbf{Standard Deviation} ($\sigma$).
Measures the variability of episode-level cumulative rewards:
\begin{equation}
\label{eq:sigma}
  \sigma = \sqrt{
    \frac{1}{E}\sum_{e=1}^{E}
    \left(\sum_{t=1}^{T} R_t^{(e)} - \mu\right)^{\!2}
  }.
\end{equation}

\textbf{Stability Coefficient} ($C_v$).
A normalized indicator of behavioral consistency, defined as
the coefficient of variation:
\begin{equation}
\label{eq:stability}
  C_v = \frac{\sigma}{|\mu|}.
\end{equation}
A lower $C_v$ indicates predictable agent behavior; a higher
value indicates volatile performance across episodes.

\textbf{Hallucination Rate} ($\eta$).
The fraction of actions classified as hallucinations, encompassing
both syntactic errors (unparseable output) and semantic errors
(logically contradictory actions), over all action decisions.
Following Definition~1 (Eq.~\ref{eq:hallucination_def}), we
adapt the hallucination taxonomy to the CAGE~4 context:
\begin{equation}
\label{eq:hallucination_rate}
  \eta = \frac{
    N_{\text{syntactic}} + N_{\text{semantic}}
  }{
    E \times T
  },
\end{equation}
where $N_{\text{syntactic}}$ counts actions that fail
schema validation (unparseable JSON, invalid action names) and
$N_{\text{semantic}}$ counts actions that pass syntax checks
but violate the IOC-expert logic $\pi_E$
(Eq.~\ref{eq:ioc_expert}) — for example, issuing \texttt{Restore}
on an uncompromised host or \texttt{Monitor} during active
compromise.

\textbf{CoSC Repair Success Rate} ($\rho$).
The fraction of hallucinated actions that the Chain-of-Self-Correction
gate successfully repaired within $K$ retries:
\begin{equation}
\label{eq:repair_rate}
  \rho = \frac{N_{\text{repaired}}}{N_{\text{syntactic}} + N_{\text{semantic}}}.
\end{equation}

\textbf{Red Impact Count} ($I_{\text{red}}$).
The total number of steps in which the red agent
successfully executed its final-stage \texttt{Impact} action,
rendering the critical OT service
unavailable~\cite{arXiv:2410.17351}:
\begin{equation}
\label{eq:red_impact}
  I_{\text{red}} = \sum_{e=1}^{E}\sum_{t=1}^{T}
    \mathbb{1}\!\bigl[\text{OT service impacted at step } t
    \text{ of episode } e\bigr].
\end{equation}
Our experiments report $I_{\text{red}}=0$ for both model
configurations, indicating complete suppression of the
adversary's strategic objective.


%=============================================================================
%  SECTION 4: METHODOLOGY — IOC Observation Enhancement
%=============================================================================
%
%  INSERT INTO: Section 4.1 — inside the Adapter Interface bullet
%  PURPOSE:     Formalize the IOC-priority encoding you mention in prose.
%               This shows exactly how Singh et al.'s observation
%               enhancement maps into your prompt structure.
%  SOURCE:      Singh et al. (2025) Section 4.2, Figure 3
%-----------------------------------------------------------------------------

% --- Equation: IOC Priority Encoding ---

Each host $v \in V$ in the agent's subnet is assigned
an IOC priority score $p(v) \in \{0,1,2,3\}$ based on
detected indicators of compromise~\cite{arXiv:2410.17351}:
\begin{equation}
\label{eq:ioc_priority}
  p(v) =
  \begin{cases}
    3, & \text{if decoy-access IOCs detected (reconnaissance)}, \\
    2, & \text{if malicious files with user-level access (exploit)}, \\
    1, & \text{if malicious files with root access (escalation)}, \\
    0, & \text{if no IOCs detected (clean)}.
  \end{cases}
\end{equation}
This priority encoding is injected into the natural language
prompt for each host, providing the LLM with the same
IOC-level granularity that yields a 42\% reward improvement
for RL agents~\cite{arXiv:2410.17351}.


%=============================================================================
%  SECTION 5: EXPERIMENTAL RESULTS — Behavioral Consistency
%=============================================================================
%
%  INSERT INTO: Section 7 — "Critical Analysis", subsection
%               "Behavioral Consistency Analysis" (already exists)
%  PURPOSE:     Formalize the improvement percentage calculation
%               that you report in Table I ("vs. Base" column).
%  SOURCE:      Standard metric definition
%-----------------------------------------------------------------------------

% --- Equation: Relative Improvement ---

The \textbf{relative improvement} over the original LLM baseline
is computed as:
\begin{equation}
\label{eq:improvement}
  \Delta = \frac{|\mu_{\text{ours}}| - |\mu_{\text{base}}|}{|\mu_{\text{base}}|}
         = \frac{|\mu_{\text{base}}| - |\mu_{\text{ours}}|}{|\mu_{\text{base}}|}
         \times 100\%,
\end{equation}
where $\mu_{\text{base}} = -2547.2$ (Castro et al.\ ALL-LLM)
and $\mu_{\text{ours}}$ is the mean reward of our agent.
For DeepSeek-R1-1.5b with $\mu=-477.5$, this yields
$\Delta = \frac{2547.2 - 477.5}{2547.2} \approx 81.2\%$.


%=============================================================================
%  SECTION 9: FUTURE RESEARCH DIRECTIONS
%=============================================================================
%
%  INSERT INTO: Section "Future Research Directions", bullet on
%               "RL-to-LLM Behavioral Distillation"
%  PURPOSE:     Show the LoRA loss equation that would be used for
%               the future work you describe.  Referencing eq:lora_loss
%               gives this future direction technical substance.
%  SOURCE:      Gao et al. (2026) Equation 2
%-----------------------------------------------------------------------------

% --- Reference only (equation already defined above as eq:lora_loss) ---
% In the Future Work bullet, add this sentence:

Specifically, the fine-tuning would minimize
the cross-entropy loss $L(w)$ (Eq.~\ref{eq:lora_loss})
over trajectories generated by the KEEP agent,
using LoRA adapters at rank~16 to fit within consumer-grade
GPU memory constraints.


%=============================================================================
%  SUMMARY: COMPLETE EQUATION REFERENCE MAP
%=============================================================================
%
%  Eq. Label              | Source Paper      | Your Section  | Purpose
%  -----------------------|-------------------|---------------|------------------
%  eq:dec_pomdp           | Singh/King        | §2.1          | Problem formulation
%  eq:gcn                 | King et al.       | §2.2          | KEEP architecture
%  eq:ppo                 | Schulman/Singh    | §2.2          | PPO training obj.
%  eq:consistency         | Hammar et al.     | §3.3          | Consistency fn λ
%  eq:abstention          | Hammar et al.     | §3.3          | Abstention policy
%  eq:gamma_calibration   | Hammar et al.     | §3.3          | Threshold tuning
%  eq:hallucination_def   | Hammar et al.     | §3.3          | Def. of hallucination
%  eq:lora_loss           | Gao et al.        | §3.4          | LoRA fine-tuning
%  eq:mc_q                | Gao et al.        | §3.4          | MC planning Q-fn
%  eq:mc_argmin           | Gao et al.        | §3.4          | Action selection
%  eq:ioc_expert          | Singh et al.      | §4.1          | CoSC expert rule
%  eq:ioc_priority        | Singh et al.      | §4.1          | IOC encoding
%  eq:joint_reward        | Castro et al.     | §4.4          | Mean reward μ
%  eq:sigma               | Standard          | §4.4          | Std. deviation σ
%  eq:stability           | Yours             | §4.4          | Stability coeff. Cv
%  eq:hallucination_rate  | Yours + Hammar    | §4.4          | Hallucination rate η
%  eq:repair_rate         | Yours             | §4.4          | CoSC repair rate ρ
%  eq:red_impact          | Singh et al.      | §4.4          | Red impact count
%  eq:improvement         | Standard          | §7            | % improvement calc.
%
%  Algorithm Label        | Source            | Your Section  | Purpose
%  -----------------------|-------------------|---------------|------------------
%  alg:cosc               | YOUR CONTRIBUTION | §4.1          | CoSC mechanism
%
%=============================================================================